\documentclass[lettersize,journal]{IEEEtran}
\IEEEoverridecommandlockouts

\usepackage{cite}
\usepackage{amsmath,amssymb,amsfonts}
\usepackage{graphicx}
\usepackage{textcomp}
\usepackage{xcolor}
\usepackage[utf8]{inputenc}
\usepackage[T1]{fontenc}

\graphicspath{{figures/}}

\def\BibTeX{{\rm B\kern-.05em{\sc i\kern-.025em b}\kern-.08em
    T\kern-.1667em\lower.7ex\hbox{E}\kern-.125emX}}

\begin{document}

\title{GNN-MARL for Network Slicing \\ (judul sementara --- ganti)}

\author{Habb,~\IEEEmembership{Student Member,~IEEE}% <-- ganti nama/afiliasi
\thanks{Manuscript received ...; revised ....}%
\thanks{Author is with the Department of ..., University, City, Country
(e-mail: email@domain.com).}%
}

\markboth{IEEE Transactions on ..., Vol.~XX, No.~X, Month~Year}%
{Habb \MakeLowercase{\textit{et al.}}: GNN-MARL for Network Slicing}

\maketitle

\begin{abstract}
TODO: abstract.
\end{abstract}

\begin{IEEEkeywords}
Network slicing, graph neural network, multi-agent reinforcement learning.
\end{IEEEkeywords}

\IEEEpeerreviewmaketitle

\section{Introduction}
\IEEEPARstart{T}{his} template sudah siap di-build. Mulai tulis di sini.
% \cite{contoh2024}

\section{Related Work}

\section{System Model}

\section{Methodology}

\section{Experiments and Results}

\section{Conclusion}

% \nocite{*} menampilkan SEMUA entri refs.bib supaya template ini langsung
% ter-compile walau belum ada \cite. Hapus baris ini begitu mulai mengutip
% referensi sungguhan dengan \cite{citekey}.
\nocite{*}
\bibliographystyle{IEEEtran}
\bibliography{refs}

\end{document}
